\subsection{Matrix Norms}

Throughout this subsection, $\mathbb{F}\in\{\mathbb{R},\mathbb{C}\}$ and $m,n,k\in\mathbb{Z}_{>0}$ whenever they occur as dimensions.

\begin{definition} \label{def:spectral-norm-and-singular-extremes} \hlt{Spectral Norm and Extreme Singular Values}\\
Let $A\in\operatorname{Mat}(m,n,\mathbb{F})$ and $q\defeq\min(m,n)$.
\begin{enumerate}[label=(\roman*)]
\setlength{\itemsep}{0pt}
\item The \hlt{spectral norm} of $A$ is $\norm{A}_2\defeq\sup\limits_{\norm{\mathbf{x}}_2=1}\norm{A\mathbf{x}}_2$.
\item If $\sigma_1\geq\cdots\geq\sigma_q\geq0$ are the singular values of $A$, set $(\forall j\in J_q)\,\sigma_j(A)\defeq\sigma_j$ and $\sigma_{\max}(A)\defeq\sigma_1$.
\item If $m=n$, set $\sigma_{\min}(A)\defeq\sigma_n$.
\end{enumerate}
\end{definition}

\begin{theorem} \label{thm:spectral-norm-largest-singular-value} \hlt{Spectral Norm is the Largest Singular Value}\\
Let $A\in\operatorname{Mat}(m,n,\mathbb{F})$ have singular values $\sigma_1\geq\cdots\geq\sigma_q\geq0$.
\begin{enumerate}[label=(\roman*)]
\setlength{\itemsep}{0pt}
\item $\norm{A}_2=\sigma_{\max}(A)$.
\item If $A=V\Sigma W^*$ is an SVD and $\mathbf{w}_1$ is first column of $W$, then $\norm{A\mathbf{w}_1}_2=\sigma_1$. Hence supremum in Definition \ref{def:spectral-norm-and-singular-extremes}(i) is attained.
\end{enumerate}
\begin{proof}
Fix an SVD $A=V\Sigma W^*$ by Theorem \ref{thm:matrix-singular-value-decomposition}, and let $\mathbf{w}_1,\ldots,\mathbf{w}_n$ be the columns of $W$.
\begin{enumerate}[label=(\roman*)]
\setlength{\itemsep}{0pt}
\item For $\mathbf{x}=\sum\limits_{j=1}^{n}c_j\mathbf{w}_j$ with $\norm{\mathbf{x}}_2=1$, $\sum\limits_{j=1}^{n}\abs{c_j}^2=1$ by Theorem \ref{thm:unitary-matrix-equivalent-conditions}(vii). Thus $\norm{A\mathbf{x}}_2^2=\sum\limits_{j=1}^{q}\sigma_j^2\abs{c_j}^2\leq\sigma_1^2$ by the SVD, so $\norm{A}_2\leq\sigma_1$.
\item The SVD has $A\mathbf{w}_1=\sigma_1\mathbf{v}_1$ with unit vectors $\mathbf{w}_1,\mathbf{v}_1$, so $\norm{A\mathbf{w}_1}_2=\sigma_1$. Hence $\norm{A}_2\geq\sigma_1$, attained at $\mathbf{w}_1$.
\end{enumerate}
\hlt{Structural Commentary}: The SVD reduces unit-vector stretch to a weighted coefficient sum.\\
\hlt{Key Insight}: The coefficient bound is sharp at $\mathbf{w}_1$.
\end{proof}
\end{theorem}

\begin{proposition} \label{prop:spectral-norm-basic-properties} \hlt{Basic Properties of the Spectral Norm}\\
Let $A,B\in\operatorname{Mat}(m,n,\mathbb{F})$, $C\in\operatorname{Mat}(n,k,\mathbb{F})$, and $c\in\mathbb{F}$.
\begin{enumerate}[label=(\roman*)]
\setlength{\itemsep}{0pt}
\item $\norm{A}_2\geq0$;
\item $\norm{A}_2=0$ if and only if $A=0$;
\item $\norm{cA}_2=\abs{c}\norm{A}_2$;
\item $\norm{A+B}_2\leq\norm{A}_2+\norm{B}_2$;
\item $\norm{AC}_2\leq\norm{A}_2\norm{C}_2$.
\end{enumerate}
\begin{proof}
If $\mathbf{y}\neq\mathbf{0}$, then $\norm{A\mathbf{y}}_2\leq\norm{A}_2\norm{\mathbf{y}}_2$ by Definition \ref{def:spectral-norm-and-singular-extremes}(i), after normalization. For $\mathbf{y}=\mathbf{0}$, the same estimate is immediate.
\begin{enumerate}[label=(\roman*)]
\setlength{\itemsep}{0pt}
\item[$(i)$] $\norm{A}_2\geq0$ by Definition \ref{def:spectral-norm-and-singular-extremes}(i).
\item[$(ii)$] ($\Rightarrow$) If $\norm{A}_2=0$, the estimate forces $(\forall\mathbf{y}\in\mathbb{F}^n)\,A\mathbf{y}=\mathbf{0}$, so $A=0$.
\item[$(ii)$] ($\Leftarrow$) If $A=0$, then $\norm{A}_2=0$ by Definition \ref{def:spectral-norm-and-singular-extremes}(i).
\item[$(iii)$] $\norm{cA}_2=\sup\limits_{\norm{\mathbf{x}}_2=1}\norm{cA\mathbf{x}}_2=\abs{c}\norm{A}_2$ by Euclidean homogeneity.
\item[$(iv)$] For $\norm{\mathbf{x}}_2=1$, $\norm{(A+B)\mathbf{x}}_2\leq\norm{A\mathbf{x}}_2+\norm{B\mathbf{x}}_2\leq\norm{A}_2+\norm{B}_2$. Thus $\norm{A+B}_2\leq\norm{A}_2+\norm{B}_2$ by the unit-sphere supremum.
\item[$(v)$] For $\norm{\mathbf{x}}_2=1$, $\norm{AC\mathbf{x}}_2\leq\norm{A}_2\norm{C\mathbf{x}}_2\leq\norm{A}_2\norm{C}_2$. Thus $\norm{AC}_2\leq\norm{A}_2\norm{C}_2$ by the unit-sphere supremum.
\end{enumerate}
\hlt{Structural Commentary}: The scaled-vector estimate feeds definiteness, homogeneity, addition, and composition.\\
\hlt{Key Insight}: Euclidean control precedes the unit-sphere supremum.
\end{proof}
\end{proposition}

\begin{definition} \label{def:operator-norm-rn-rm} \hlt{Operator Norm and Induced Metric on $\LL(\R^n,\R^m)$}\\
For $A,B\in\LL(\R^n,\R^m)$, define $\normop{A}\defeq\sup\limits_{\norm{\h}_2\leq1}\norm{A\h}_2$ and $d(A,B)\defeq\normop{A-B}$.\\
Under standard-basis identification, $\normop{A}=\norm{[A]}_2$ by scaling; hence $(\forall\h\in\R^n)\,\norm{A\h}_2\leq\normop{A}\norm{\h}_2$.\\
The norm laws for sums and compatible compositions are Proposition \ref{prop:spectral-norm-basic-properties}.
\end{definition}

\begin{proposition} \label{prop:frobenius-and-spectral-norm-comparison} \hlt{Frobenius and Spectral Norm Comparison}\\
Let $A\in\operatorname{Mat}(m,n,\mathbb{F})$ have singular values $\sigma_1\geq\cdots\geq\sigma_q\geq0$.
\begin{enumerate}[label=(\roman*)]
\setlength{\itemsep}{0pt}
\item $\norm{A}_F^2=\sum\limits_{j=1}^{q}\sigma_j^2$.
\item $\norm{A}_2\leq\norm{A}_F$.
\item If $A\neq0$, then $\norm{A}_2=\norm{A}_F$ if and only if $\operatorname{rank}(A)=1$.
\end{enumerate}
\begin{proof}
\begin{enumerate}[label=(\roman*)]
\setlength{\itemsep}{0pt}
\item[$(i)$] The identity holds by Proposition \ref{prop:basic-singular-value-identities}(ii).
\item[$(ii)$] $\norm{A}_2^2=\sigma_1^2\leq\sum\limits_{j=1}^{q}\sigma_j^2=\norm{A}_F^2$ by Theorem \ref{thm:spectral-norm-largest-singular-value}(i) and part (i).
\item[$(iii)$] ($\Rightarrow$) If $A\neq0$ and $\norm{A}_2=\norm{A}_F$, then $\sum\limits_{j=2}^{q}\sigma_j^2=\norm{A}_F^2-\norm{A}_2^2=0$ by parts (i)--(ii). Thus $\operatorname{rank}(A)=1$ by Definition \ref{def:matrix-singular-values}.
\item[$(iii)$] ($\Leftarrow$) If $\operatorname{rank}(A)=1$, then $\sigma_2=\cdots=\sigma_q=0$ by Definition \ref{def:matrix-singular-values}; hence $\norm{A}_2=\norm{A}_F$ by parts (i)--(ii).
\end{enumerate}
\hlt{Structural Commentary}: Frobenius aggregates all squared singular values; spectral isolates their maximum.\\
\hlt{Key Insight}: Equality excludes every singular direction except one.
\end{proof}
\end{proposition}

\begin{theorem} \label{thm:frobenius-and-spectral-submultiplicativity} \hlt{Submultiplicativity of Frobenius and Spectral Norms}\\
Let $A\in\operatorname{Mat}(m,n,\mathbb{F})$ and $B\in\operatorname{Mat}(n,k,\mathbb{F})$.
\begin{enumerate}[label=(\roman*)]
\setlength{\itemsep}{0pt}
\item $\norm{AB}_F\leq\norm{A}_2\norm{B}_F\leq\norm{A}_F\norm{B}_F$.
\item $\norm{AB}_2\leq\norm{A}_2\norm{B}_2$.
\end{enumerate}
\begin{proof}
\begin{enumerate}[label=(\roman*)]
\setlength{\itemsep}{0pt}
\item Write $B=[\mathbf{b}_1\ \cdots\ \mathbf{b}_k]$. Then $(\forall j\in J_k)\,\norm{A\mathbf{b}_j}_2\leq\norm{A}_2\norm{\mathbf{b}_j}_2$ by Definition \ref{def:spectral-norm-and-singular-extremes}(i) and scaling.\\
Hence $\norm{AB}_F^2=\sum\limits_{j=1}^{k}\norm{A\mathbf{b}_j}_2^2\leq\norm{A}_2^2\sum\limits_{j=1}^{k}\norm{\mathbf{b}_j}_2^2=\norm{A}_2^2\norm{B}_F^2$; $\norm{AB}_F\leq\norm{A}_2\norm{B}_F\leq\norm{A}_F\norm{B}_F$ by Proposition \ref{prop:frobenius-and-spectral-norm-comparison}(ii).
\item The spectral-norm inequality holds by Proposition \ref{prop:spectral-norm-basic-properties}(v).
\end{enumerate}
\hlt{Structural Commentary}: Columnwise spectral bounds sum to Frobenius control; norm comparison closes the chain.\\
\hlt{Key Insight}: Each column of $AB$ inherits the same factor $\norm{A}_2$.
\end{proof}
\end{theorem}

\begin{proposition} \label{prop:unitary-invariance-of-singular-values} \hlt{Unitary Invariance of Singular Values}\\
Let $A\in\operatorname{Mat}(m,n,\mathbb{F})$, and let $U\in\operatorname{Mat}(m,m,\mathbb{F})$ and $V\in\operatorname{Mat}(n,n,\mathbb{F})$ be unitary.
\begin{enumerate}[label=(\roman*)]
\setlength{\itemsep}{0pt}
\item $A$ and $UAV$ have the same singular values.
\item $\norm{UAV}_2=\norm{A}_2$.
\item $\norm{UAV}_F=\norm{A}_F$.
\end{enumerate}
\begin{proof}
\begin{enumerate}[label=(\roman*)]
\setlength{\itemsep}{0pt}
\item Fix an SVD $A=X\Sigma Y^*$ by Theorem \ref{thm:matrix-singular-value-decomposition}.\\
Since $UX$ and $V^*Y$ are unitary, $UAV=(UX)\Sigma(V^*Y)^*$ is an SVD with the same diagonal $\Sigma$.
\item $\norm{UAV}_2=\norm{A}_2$ by part (i) and Theorem \ref{thm:spectral-norm-largest-singular-value}(i).
\item $\norm{UAV}_F=\norm{A}_F$ by part (i) and Proposition \ref{prop:basic-singular-value-identities}(ii).
\end{enumerate}
\hlt{Structural Commentary}: Multiplying the two SVD bases absorbs the exterior unitary factors.\\
\hlt{Key Insight}: The diagonal singular factor $\Sigma$ remains unchanged.
\end{proof}
\end{proposition}

\begin{counterexample} \label{cex:similarity-need-not-preserve-singular-values} \hlt{Similarity Need Not Preserve Singular Values}\\
The claim that arbitrary similarity preserves singular values fails in $\operatorname{Mat}(2,2,\mathbb{R})$. Set $A=\left(\begin{smallmatrix}1&1\\0&1\end{smallmatrix}\right)$, $B=\left(\begin{smallmatrix}1&2\\0&1\end{smallmatrix}\right)$, and $D=\operatorname{diag}(2,1)$.\\
Then $B=DAD^{-1}$, so $A$ and $B$ are similar. Their Gram matrices are $A^*A=\left(\begin{smallmatrix}1&1\\1&2\end{smallmatrix}\right)$ and $B^*B=\left(\begin{smallmatrix}1&2\\2&5\end{smallmatrix}\right)$.\\
Both $A$ and $B$ are invertible, so all their Gram eigenvalues are positive. The Gram matrices have traces $3$ and $6$.\\
Hence their spectra, and therefore the singular values of $A$ and $B$, differ by Proposition \ref{prop:basic-singular-value-identities}(i).\\
Thus Proposition \ref{prop:unitary-invariance-of-singular-values} cannot be extended from unitary left-right equivalence to arbitrary similarity.
\end{counterexample}

\begin{theorem} \label{thm:schur-inequality-defect-normality} \hlt{Schur's Inequality and Normality}\\
Let $A\in\operatorname{Mat}(n,n,\mathbb{C})$ have eigenvalues $\lambda_1,\ldots,\lambda_n$, counted with algebraic multiplicity.
\begin{enumerate}[label=(\roman*)]
\setlength{\itemsep}{0pt}
\item $\sum\limits_{j=1}^{n}\abs{\lambda_j}^2\leq\norm{A}_F^2$.
\item Equality holds in (i) if and only if $A$ is normal.
\end{enumerate}
\begin{proof}
Fix a unitary $U$ such that $T\defeq U^*AU=[t_{ij}]$ is upper triangular by Theorem \ref{thm:simultaneous-schur-triangularization-commuting-family}(i) with $\mathcal{F}=\{A\}$.\\
Relabel the eigenvalues so that $(\forall j\in J_n)\,t_{jj}=\lambda_j$ by Proposition \ref{prop:upper-triangular-characteristic-polynomial}(iii).
\begin{enumerate}[label=(\roman*)]
\setlength{\itemsep}{0pt}
\item[$(i)$] $\norm{A}_F^2=\norm{T}_F^2=\sum\limits_{i,j=1}^{n}\abs{t_{ij}}^2$ by Propositions \ref{prop:unitary-invariance-of-singular-values}(iii) and \ref{prop:basic-singular-value-identities}(ii), with direct matrix multiplication.\\
Thus $\norm{A}_F^2=\sum\limits_{j=1}^{n}\abs{\lambda_j}^2+\sum\limits_{i<j}\abs{t_{ij}}^2$ by triangularity, and (i) holds by nonnegativity of the final sum.
\item[$(ii)$] ($\Rightarrow$) Equality forces $\sum\limits_{i<j}\abs{t_{ij}}^2=0$ by part (i), so $T$ is diagonal and $A=UTU^*$ is normal.
\item[$(ii)$] ($\Leftarrow$) If $A$ is normal, then $T^*T=U^*A^*AU=U^*AA^*U=TT^*$, so $T$ is diagonal by Lemma \ref{lem:upper-triangular-normal-matrix-is-diagonal}. Equality follows by part (i).
\end{enumerate}
\hlt{Structural Commentary}: Unitary Schur form partitions Frobenius energy into diagonal spectral and strict-upper terms.\\
\hlt{Key Insight}: Equality is equivalent to vanishing strict-upper energy.
\end{proof}
\end{theorem}

\begin{definition} \label{def:defect-from-normality} \hlt{Defect from Normality}\\
Let $A\in\operatorname{Mat}(n,n,\mathbb{C})$ have eigenvalues $\lambda_1,\ldots,\lambda_n$ with algebraic multiplicity. Its \hlt{defect from normality} is $\Delta(A)\defeq\norm{A}_F^2-\sum\limits_{j=1}^{n}\abs{\lambda_j}^2$.\\
Then $\Delta(A)\geq0$, and $\Delta(A)=0$ if and only if $A$ is normal, by Theorem \ref{thm:schur-inequality-defect-normality}(i)--(ii).
\end{definition}

\begin{theorem} \label{thm:normality-transfer-ab-ba} \hlt{Normality Transfers from $AB$ to $BA$}\\
Let $A,B\in\operatorname{Mat}(n,n,\C)$. If $A$, $B$, and $AB$ are normal, then $BA$ is normal.
\begin{proof} $AB, BA$ have same eigenvalue multiset by Theorem \ref{thm:ab-ba-nonzero-eigenvalues}(i) with equal sizes. Set $M\defeq A^*A=AA^*$, $N\defeq B^*B=BB^*$.\\
Then $\norm{BA}_F^2=\operatorname{tr}(A^*NA)=\operatorname{tr}(MN)$ and $\norm{AB}_F^2=\operatorname{tr}(B^*MB)=\operatorname{tr}(NM)=\operatorname{tr}(MN)$ by Propositions \ref{prop:basic-singular-value-identities}(ii) and \ref{prop:trace-cyclic-products}(i).\\
Hence $\Delta(BA)=\Delta(AB)$ by Definition \ref{def:defect-from-normality}, $\Delta(AB)=0$ by Theorem \ref{thm:schur-inequality-defect-normality}(ii); $\Delta(BA)=0$, so $BA$ is normal by Theorem \ref{thm:schur-inequality-defect-normality}(ii).\\
\hlt{Structural Commentary}: Common spectral data and a cyclic trace identity transfer the Schur defect from $AB$ to $BA$.\\
\hlt{Key Insight}: Normal factors force $\norm{BA}_F=\norm{AB}_F$.
\end{proof}
\end{theorem}

\begin{theorem} \label{thm:eigenvalues-versus-singular-values} \hlt{Eigenvalues Versus Singular Values}\\
Let $A\in\operatorname{Mat}(n,n,\mathbb{C})$ have eigenvalues $\lambda_1,\ldots,\lambda_n$ ordered so that $\abs{\lambda_1}\geq\cdots\geq\abs{\lambda_n}$.\\
Let $\sigma_1\geq\cdots\geq\sigma_n\geq0$ be its singular values.
\begin{enumerate}[label=(\roman*)]
\setlength{\itemsep}{0pt}
\item $\abs{\det A}=\prod\limits_{j=1}^{n}\abs{\lambda_j}=\prod\limits_{j=1}^{n}\sigma_j$;
\item $\rho(A)\leq\norm{A}_2$;
\item $\sigma_n\leq\abs{\lambda_n}$;
\item $A$ is normal if and only if $(\forall j\in J_n)\,\abs{\lambda_j}=\sigma_j$;
\item $A$ is positive semidefinite if and only if $(\forall j\in J_n)\,\lambda_j=\sigma_j$.
\end{enumerate}
\begin{proof}
\begin{enumerate}[label=(\roman*)]
\setlength{\itemsep}{0pt}
\item[$(i)$] $\abs{\det A}=\prod\limits_{j=1}^{n}\abs{\lambda_j}$ by Proposition \ref{prop:trace-determinant-from-eigenvalues}(ii).\\
For an SVD $A=V\Sigma W^*$, $\abs{\det A}=\abs{\det\Sigma}=\prod\limits_{j=1}^{n}\sigma_j$ by Proposition \ref{prop:isometry-determinant-absolute-one}(i) and Theorem \ref{thm:determinant-basic-laws}(v)--(vi).
\item[$(ii)$] If $A\mathbf{x}=\lambda\mathbf{x}$ and $\norm{\mathbf{x}}_2=1$, then $\abs{\lambda}=\norm{A\mathbf{x}}_2\leq\norm{A}_2$ by Definition \ref{def:spectral-norm-and-singular-extremes}(i); hence (ii) by maximizing over $\lambda$.
\item[$(iii)$] If $A$ is singular, $\abs{\lambda_n}=\sigma_n=0$ by Corollaries \ref{cor:matrix-invertible-iff-zero-not-spectrum} and \ref{cor:square-invertibility-singular-value-dictionary}.\\
If $A$ is invertible, then $\rho(A^{-1})=\abs{\lambda_n}^{-1}$ by Corollary \ref{cor:eigenvalues-of-inverse-matrix}. For an SVD $A=V\Sigma W^*$, $A^{-1}=W\Sigma^{-1}V^*$, so $\norm{A^{-1}}_2=\sigma_n^{-1}$ by Theorem \ref{thm:spectral-norm-largest-singular-value}(i), after reversing the diagonal order. Thus $\sigma_n\leq\abs{\lambda_n}$ by part (ii) for $A^{-1}$.
\item[$(iv)$] ($\Rightarrow$) After reordering, fix $A=U\operatorname{diag}(\lambda_1,\ldots,\lambda_n)U^*$ by Theorem \ref{thm:complex-spectral-theorem}(i)$\Rightarrow$(ii).\\
Then $(\forall j\in J_n)\,\sigma_j=\abs{\lambda_j}$ by Proposition \ref{prop:unitary-invariance-of-singular-values}(i) and Definition \ref{def:matrix-singular-values}.
\item[$(iv)$] ($\Leftarrow$) If $(\forall j\in J_n)\,\abs{\lambda_j}=\sigma_j$, then $\sum\limits_{j=1}^{n}\abs{\lambda_j}^2=\sum\limits_{j=1}^{n}\sigma_j^2=\norm{A}_F^2$ by Proposition \ref{prop:basic-singular-value-identities}(ii). Thus $A$ is normal by Theorem \ref{thm:schur-inequality-defect-normality}(ii).
\item[$(v)$] ($\Rightarrow$) If $A\succeq0$, then $A$ is Hermitian with nonnegative eigenvalues by Theorem \ref{thm:positive-semidefinite-matrix-equivalences}(i)$\Rightarrow$(ii).\\
Hence $A$ is normal and $(\forall j\in J_n)\,\lambda_j=\sigma_j$ by part (iv).
\item[$(v)$] ($\Leftarrow$) If $(\forall j\in J_n)\,\lambda_j=\sigma_j$, then $A$ is normal by part (iv) and its eigenvalues are nonnegative.\\
By Theorem \ref{thm:complex-spectral-theorem}(i)$\Rightarrow$(ii), $A=U\operatorname{diag}(\lambda_1,\ldots,\lambda_n)U^*$ for a unitary $U$, so $A=A^*$ and $A\succeq0$ by Theorem \ref{thm:positive-semidefinite-matrix-equivalences}(ii)$\Rightarrow$(i).
\end{enumerate}
\hlt{Structural Commentary}: Determinants, extremal norms, and Schur energy compare eigenvalues with singular values.\\
\hlt{Key Insight}: Equality of squared spectral masses detects normality.
\end{proof}
\end{theorem}

\begin{theorem} \label{thm:trace-bounded-by-singular-values} \hlt{Trace Bound by Singular Values}\\
Let $A=[a_{ij}]\in\operatorname{Mat}(n,n,\mathbb{C})$ have eigenvalues $\lambda_1,\ldots,\lambda_n$, counted with algebraic multiplicity. Let $\sigma_1,\ldots,\sigma_n$ be its singular values.
\begin{enumerate}[label=(\roman*)]
\setlength{\itemsep}{0pt}
\item $\abs{\operatorname{tr}A}\leq\sum\limits_{j=1}^{n}\sigma_j$;
\item $\sum\limits_{j=1}^{n}\abs{\lambda_j}\leq\sum\limits_{j=1}^{n}\sigma_j$;
\item $\sum\limits_{j=1}^{n}\abs{a_{jj}}\leq\sum\limits_{j=1}^{n}\sigma_j$.
\end{enumerate}
\begin{proof}
Set $\omega:\mathbb{C}\to\mathbb{C}$ by $\omega(0)\defeq1$ and $\omega(z)\defeq\overline{z}/\abs{z}$ for $z\neq0$; then $(\forall z\in\mathbb{C})\,\abs{\omega(z)}=1$ and $\omega(z)z=\abs{z}$.
\begin{enumerate}[label=(\roman*)]
\setlength{\itemsep}{0pt}
\item Fix an SVD $A=V\Sigma W^*$ by Theorem \ref{thm:matrix-singular-value-decomposition}, and write $\mathbf{v}_j,\mathbf{w}_j$ for the columns of $V,W$.\\
Then $\operatorname{tr}A=\sum\limits_{j=1}^{n}\sigma_j\mathbf{w}_j^*\mathbf{v}_j$ by Proposition \ref{prop:trace-cyclicity-two-factors}, while $(\forall j\in J_n)\,\abs{\mathbf{w}_j^*\mathbf{v}_j}\leq1$ by Theorem \ref{thm:cauchy-schwarz-inequality}(i); $\abs{\operatorname{tr}A}\leq\sum\limits_{j=1}^{n}\sigma_j$ by the triangle inequality.
\item Fix a Schur form $T\defeq U^*AU=[t_{ij}]$ by Theorem \ref{thm:simultaneous-schur-triangularization-commuting-family}(i) with $\mathcal{F}=\{A\}$.\\
Relabel so that $(\forall j\in J_n)\,t_{jj}=\lambda_j$ by Proposition \ref{prop:upper-triangular-characteristic-polynomial}(iii). For $D_\lambda\defeq\operatorname{diag}(\omega(\lambda_1),\ldots,\omega(\lambda_n))$, $\operatorname{tr}(D_\lambda T)=\sum\limits_{j=1}^{n}\abs{\lambda_j}$.\\
The matrices $D_\lambda T$ and $A$ have the same singular values by Proposition \ref{prop:unitary-invariance-of-singular-values}(i), so (ii) holds by part (i) for $D_\lambda T$.
\item For $D_A\defeq\operatorname{diag}(\omega(a_{11}),\ldots,\omega(a_{nn}))$, $\operatorname{tr}(D_AA)=\sum\limits_{j=1}^{n}\abs{a_{jj}}$.\\
The matrices $D_AA$ and $A$ have the same singular values by Proposition \ref{prop:unitary-invariance-of-singular-values}(i), so (iii) holds by part (i) for $D_AA$.
\end{enumerate}
\hlt{Structural Commentary}: SVD controls trace; diagonal phase rotations transfer that bound to Schur and coordinate diagonals.\\
\hlt{Key Insight}: Choose phases so every selected diagonal contribution is nonnegative.
\end{proof}
\end{theorem}
